% Generated draft: requires mathematical and source review.
\documentclass[11pt]{article}
\usepackage[margin=25mm]{geometry}
\usepackage{fontspec}
\setmainfont{Latin Modern Roman}
\usepackage{amsmath,amssymb,mathtools,unicode-math}
\setmathfont{Latin Modern Math}
\usepackage{enumitem,needspace,xurl,hyperref}
\setlength{\parindent}{0pt}
\setlength{\parskip}{5pt}
\setlength{\emergencystretch}{3em}
\newcommand{\N}{\mathbb{N}}
\newcommand{\Z}{\mathbb{Z}}
\newcommand{\Q}{\mathbb{Q}}
\newcommand{\R}{\mathbb{R}}
\newcommand{\C}{\mathbb{C}}
\newcommand{\F}{\mathbb{F}}
\newcommand{\A}{\mathbb{A}}
\newcommand{\PP}{\mathbb{P}}
\newcommand{\E}{\mathbb{E}}
\newcommand{\eps}{\varepsilon}
\newcommand{\dd}{\,\mathrm{d}}
\begin{document}
\section*{UnsolvedMath ID: 30002616}
\textbf{Source problem code:} catalogue entry
\subsection{Definitions and mathematical statement}
Let $E/\Q$ be an elliptic curve, that is, a smooth projective curve of genus one with a specified rational point. An Artin representation $\rho$ is a continuous representation of $\operatorname{Gal}(\overline{\Q}/\Q)$ with finite image on a finite-dimensional vector space over a characteristic-zero coefficient field. The supplied record does not specify that coefficient field or the additional restrictions on $E$ and $\rho$ needed for the conjecture it summarizes.

Write $L(E,\rho,s)$ for the Artin-twisted $L$-function referred to in the record. For a field extension $H/\Q$, $E(H)$ denotes the group of $H$-rational points of $E$. The record uses $E(H)_\rho$ for a $\rho$-component, but does not specify $H$, its relation to $\rho$, or the convention and coefficient field defining this component. It also does not specify the prime $p$ or its required local hypotheses.

Denote the generalized Kato class mentioned in the record by $\kappa$, and denote its unspecified notion of equivariant analytic rank by $r_\rho$. These are placeholders: the construction and ambient cohomology group of $\kappa$, the comparison map needed to interpret membership in a rational-point space, and the definition of $r_\rho$ are absent. Subject to the missing construction hypotheses and conventions, the supplied question has the schematic form
\[
L(E,\rho,1)=0
\quad\Longrightarrow\quad
\begin{cases}
\kappa\in E(H)_\rho\otimes\Q_p,\\
\kappa=0 & \text{if }r_\rho>2.
\end{cases}
\]
When $r_\rho=2$, the record additionally asks whether $\kappa$ is given by a proposed $p$-adic-period formula. That formula, including its periods, normalizations, and choices, is not supplied and cannot be stated precisely from this record.

The supplied curation describes this as a truncated Darmon--Rotger conjecture concerning rank two elliptic curves and vanishing beyond rank two. Accordingly, the display records the intended questions only; it is not a complete conjecture quantified over arbitrary elliptic curves and Artin representations. The missing hypotheses, rank convention, class construction, and rank-two formula require recovery from the cited sources before a precise mathematical formulation can be asserted.
\subsection{Short English statement}
Assuming the Artin-twisted central value vanishes, the supplied question predicts that a generalized Kato class comes from rational points, vanishes when the stated equivariant analytic rank exceeds two, and satisfies a proposed $p$-adic-period formula in rank two. Essential hypotheses, definitions, and the formula itself are missing, so this is a statement-only schematic draft requiring human mathematical review.
\subsection{Sources}
\begin{itemize}
\item Local UnsolvedMath registry: supplied record with permanent ID $30002616$, title ``Generalized Kato Classes at Vanishing Artin-Twisted $L$-Values,'' exact source problem code \texttt{catalogue entry}, in the Oberwolfach Reports Open Problems collection. The registry gives proposed year 2014 and supplies no individual proposer.
\item Supplied source citation: \emph{Algebraische Zahlentheorie (2015)}, Oberwolfach Reports, pp.~1747--1800, DOI \url{https://doi.org/10.4171/owr/2014/32}. The background describes Oberwolfach Reports as documenting workshops and their open-problem discussions. The citation's 2015 designation and its 2014 DOI/year metadata are retained without reconciliation.
\item Supplied literature source: Darmon--Rotger, ``elliptic curves and generalized Kato classes,'' \url{https://doi.org/10.1186/s40687-016-0074-9}. The supplied curation attributes the conjectural relation to Darmon--Rotger; this is the supplied reference label, not an independently verified bibliographic title.
\item Supplied literature source: ``original Oberwolfach report,'' \url{https://doi.org/10.4171/owr/2014/32}.
\item Literature assessment dated exactly \texttt{2026-08-21}: the supplied curation labels the problem open, describes the record as truncating a Darmon--Rotger conjecture, and says the primary paper formulates the expected relation without a general proof. These are supplied dated curation claims, not independently verified conclusions. The background explicitly requests independent verification of time-sensitive status and cited claims.
\item The registry labels the statement \texttt{corrected\_verified} and says its mathematical content was preserved while notation, typography, or omitted definitions were repaired for publication. This supplied editorial label does not establish completeness or mathematical correctness of the present formulation.
\end{itemize}
\end{document}
